\subsection*{一步一步建立直观}

\paragraph{1. 先把等式当作运动约束}

令右图保持单位半径。圆方程 $x^2+y^2=1$ 要求点的两个坐标协同变化：一个变大时，另一个必须做出补偿。沿圆运动的小位移 $(dx,dy)$ 在一阶上满足 $2x\,dx+2y\,dy=0$，这比先求出上下两个平方根分支更直接。



\paragraph{2. 梯度告诉我们不能朝哪里走}

约束的梯度 $(2x,2y)$ 沿半径向外，切向位移与它垂直，所以一阶变化为零。转到圆的左端点，切线竖直；此时不能把附近的整段圆当成一个可微的 $y(x)$。隐式关系仍然有意义，某种坐标表达却可能失效。



\paragraph{3. 把可除的条件写出来}

从约束式得到 $dy/dx=-x/y$ 时，必须假设 $y\ne0$。回到圆顶，分母非零，斜率为零；回到左右端点，分母为零，提示应换用 $x(y)$ 或参数表示。公式里的限制并非技术障碍，而是图形姿态在坐标中的真实反映。


